\mysection{Evaluation}
\label{sec:eval}

We implemented a prototype of \tool with $\sim$10,400 lines of Python and 774 lines of C++ code. As the serving backend, we used SGLang~\cite{sglang@arxiv23}, a widely adopted open-source inference engine, and extended it with our elastic memory manager library.
The library is built on CUDA VMM APIs~\cite{cuda-vmm} and exposes standard KV cache allocation APIs via Python bindings, requiring only 22 lines of changes to integrate with SGLang.

On the frontend, we used a Redis queue~\cite{redis-list} to cache incoming requests from all clients. \tool's local scheduler dispatches these requests to the serving engines of corresponding models based on Algorithm~\ref{alg:local}.
For tensor-parallel models, the GPU-local scheduler runs only on the first rank, and the resulting scheduling decisions are broadcast to all other ranks to ensure consistency.
\tool's global scheduler operates as a separate Python process, collecting execution metrics from each engine\textemdash such as request rates and queue status. It makes scheduling decisions (\eg, model evictions and activations) and communicates them to the engines using ZeroMQ~\cite{zmq}.

\mysubsection{Experimental Setup}

\MyPara{Testbed.}
We conducted our experiments on a cluster of four nodes, each equipped with eight NVIDIA H100-80G GPUs interconnected via 600GB/s NVLink. These nodes communicate through a 100Gbps Ethernet network. Each node features two 52-core Intel Xeon Platinum 8480+ CPUs, 1.7 TB of DRAM, and a PCIe Gen5 x16 interface. All nodes run Ubuntu 22.04 with CUDA Toolkit 12.4.

\MyPara{Baselines.}
We compared \tool against four baselines that share GPUs to serve multi-LLMs: (1) \textit{Static partition (S-Partition)}; (2) \textit{MuxServe++}; (3) \textit{QLM~\cite{qlm}}; and (4) ServerlessLLM~\cite{fu2024serverlessllm}.
Note that the original MuxServe is built on vLLM and supports only Llama-2 models. We ported it to SGLang and generalized its memory mechanism with our kvcached to support heterogeneous models, referred to as \muxp.
We evaluated the performance of MuxServe++ and MuxServe using three Llama-3.1-8B models under different request rates over a 10-minute period: 199 requests/min, 262 requests/min, and 22 requests/min. All experiments were conducted under the same and consistent conditions. The results are shown in Table~\ref{tab:muxserve-comparison-slim}.
As we can see, MuxServe++ achieves comparable or even better performance.

\begin{table}[t]
\centering
\begin{adjustbox}{max width=\linewidth}
\begin{tabular}{lcccc}
\toprule
 & Mean E2E (s) & P95 E2E (s) & Req Tput (r/s) & Token Tput (t/s) \\
\midrule
MuxServe    & 7.40 & 18.85 & 7.97 & 3363.53 \\
MuxServe++  & 5.25 & 12.09 & 7.98 & 3353.09 \\
\midrule
 & Mean TTFT (s) & P95 TTFT (s) & Mean TPOT (ms) & P95 TPOT (ms) \\
MuxServe    & 1.47 & 11.04 & 21.21 & 31.95 \\
MuxServe++  & 0.089 & 0.320 & 18.82 & 33.97 \\
\bottomrule
\end{tabular}
\end{adjustbox}
\caption{Performance comparison of MuxServe and MuxServe++}
\label{tab:muxserve-comparison-slim}
\end{table}

Note that general GPU-sharing systems target application-agnostic workloads and cannot be directly compared: they focus on single-GPU sharing rather than cluster-level GPU sharing, and are not designed to optimize LLM serving, which requires both TTFT and TPOT adherence (\S\ref{sec:rw}).

\MyPara{Traces and models.} 
We used two real-world traces, Hyperbolic~\cite{hyperbolic} and Arena-Chat~\cite{chatbotarena@arxiv24}. For each trace, we sampled a representative set of models, including both popular and long-tail models, ensuring their workload characteristics (\eg, popularity distribution and idle patterns) align with the observations in \S\ref{sec:motivation}.
To simulate various scenarios, we scaled the traces by multiplying the number of requests by a factor of $N$, increasing the load while preserving the original traffic patterns---a common way used in prior work~\cite{alpaserve@osdi23, muxserve}.
In total, we evaluated 58 LLMs as detailed in Table~\ref{tab:models}. The large-scale experiments (\S\ref{sec:large-scale}) use all models, while other evaluations (\S\ref{sec:e2e}---\S\ref{sec:ablation}) select subsets tailored to specific goals.

\begin{table}[t]
\centering
\footnotesize
\begin{tabular}{lcccc}
\toprule
\textbf{Model size} & 1B-3B & 4B-8B & 9B-30B & 31B-70B \\
\midrule
\#LLMs & 43 & 8 & 3 & 4 \\
\bottomrule
\end{tabular}
\vspace{-2mm}
\caption{Models used in our evaluation.}
\vspace{-1mm}
\label{tab:models}
\end{table}

\begin{figure*}[th]
  \centering
  \includegraphics[width=0.8\textwidth]{figs/legend-fixed.pdf}
  \vspace{-2mm}
  
  \begin{subfigure}[t]{0.44\textwidth}
    \includegraphics[width=\linewidth]{figs/e2e/e2e-hyper-rate-scale-fixed.pdf}
  \end{subfigure}
  \begin{subfigure}[t]{0.44\textwidth}
    \includegraphics[width=\linewidth]{figs/e2e/e2e-arena-rate-scale-fixed.pdf}
  \end{subfigure}
  \vspace{-1.8mm}

  \begin{subfigure}[t]{0.44\textwidth}
    \includegraphics[width=\linewidth]{figs/e2e/e2e-hyper-slo-scale-fixed.pdf}
  \end{subfigure}
  \begin{subfigure}[t]{0.44\textwidth}
    \includegraphics[width=\linewidth]{figs/e2e/e2e-arena-slo-scale-fixed.pdf}
  \end{subfigure}
  \vspace{-1.8mm}

  \begin{subfigure}[t]{0.44\textwidth}
    \includegraphics[width=\linewidth]{figs/e2e/e2e-hyper-device-scale-fixed.pdf}
  \end{subfigure}
  \begin{subfigure}[t]{0.44\textwidth}
    \includegraphics[width=\linewidth]{figs/e2e/e2e-arena-device-scale-fixed.pdf}
  \end{subfigure}
  \vspace{-3mm}
  \caption{End-to-end performance comparison on SLO attainments under varying scales of rates, SLOs, and numbers of available GPUs. The dotted vertical lines mark where the system reaches 99\% TTFT or TPOT attainment.}
  \label{fig:slo-comparison}
  \vspace{-1em}
\end{figure*}
\MyPara{Metrics.} 
Our primary performance metrics are TTFT and TPOT attainment.
To establish SLOs for each model, we first ran its workload on dedicated GPUs to measure its P95 TTFT and TPOT latencies. This process produced TTFT SLOs ranging from 0.04s to 0.13s and TPOT SLOs from 5.2ms to 50.9ms.
We then scaled these base values by a factor to evaluate system performance under varying latency requirements, following an approach consistent with prior work~\cite{alpaserve@osdi23, muxserve, qin2025mooncake}.
We also reported aggregated throughput. To account for model idle periods, throughput considers the actual time (excluding idle time) spent serving them.

\mysubsection{End-to-End Performance} 
\label{sec:e2e}

We start with the end-to-end performance of \tool under varying request rates, SLO requirements, and GPU supplies.

\MyPara{SLO attainment vs. request rate.}
We first evaluated various inference loads by serving eight models on two shared GPUs.
As shown in Figure~\ref{fig:slo-comparison} (first row), \tool consistently outperforms all baselines by maintaining a significantly higher TTFT SLO attainment.
On the Hyperbolic trace, \tool supports up to $2.3\times$ and $3.5\times$ more requests than \muxp and static partitioning, respectively, while still achieving 99\% SLO attainment. On the Arena-Chat trace, it handles over $3\times$ more requests than all baselines. The gains stem from KVPR-driven placement, which keeps per-GPU memory pressure low so active models retain headroom under bursty arrivals.
\muxp's TTFT SLO attainment drops quickly with higher request rates because it cannot evict idle models or relocate models across GPUs, leading to memory contention and degraded performance.

QLM time-shares GPUs by packing pending requests into groups and dispatching each group to a GPU under an SLO-aware stochastic policy, swapping models and preempting unfinished requests whenever consecutive groups target different models. The frequent swapping makes QLM's TTFT attainment worse than static partitioning.

ServerlessLLM serves models in a time-sharing manner for serverless workloads: inactive models are unloaded, and newly arriving requests wait until the scheduler reactivates the required model on the server with the fastest startup, selected based on checkpoint locality. Because the full cold-start process still dominates request latency, ServerlessLLM achieves the worst TTFT SLO attainment.

\tool maintains high TPOT attainment through its demand-aware scheduler, which balances workloads to reduce memory contention. Although the scheduler explicitly targets TTFT, TPOT benefits as a side effect: TPOT is degraded mainly by oversized batches and memory-induced preemptions, both of which subside once KVPR-balanced placement caps per-GPU memory contention. QLM achieves lower TPOT because it over-batches under high load, increasing per-iteration latency, while ServerlessLLM performs even worse by allowing unbounded batch sizes that further raise P99 TPOT. Both \muxp and static partitioning experience severe memory contention at high request rates, causing frequent preemptions that substantially degrade TPOT.

\MyPara{SLO attainment vs. SLO requirements.} 
The middle row of Figure~\ref{fig:slo-comparison} shows the attainment across different SLO targets.
As SLO scales up, \tool quickly achieves 99\% TTFT and TPOT attainment. 
In contrast, no baseline achieves 99\% TTFT attainment even on the largest SLO scale in this experiment, and their attainment rates do not improve significantly as the scale increases. 
Among the baselines, \muxp achieves the best TTFT performance, reaching 84.79\% and 67.22\% attainment on the largest SLO scale for the two traces. The gap to \tool is primarily due to their inflexibility to adapt sharing policies dynamically. 
The TPOT attainment of all systems increases rapidly as the SLO scale grows 
because TPOT is less sensitive to memory contention. 
We also observe that the Hyperbolic trace requires higher SLO scales due to more bursty and heavier request patterns.

\MyPara{SLO attainment vs. available GPUs.}
Finally, we evaluate performance when provisioning more GPUs. We selected 18 models from Table~\ref{tab:models}, representing a mix of popular and tail models with diverse load variability.
To fully test the flexibility of our scheduling strategy, we included models of varying sizes from 1B to 8B, all of which fit within a single 80GB GPU. This setup enables a wide range of model-sharing combinations across GPUs.
As shown in Figure~\ref{fig:slo-comparison} (last row), \tool achieves 99\% TTFT and TPOT attainment using only four and five GPUs on the two traces, respectively, demonstrating its effectiveness in improving cost-efficiency while maintaining performance.
In comparison, all baselines fail to reach 99\% TTFT attainment even with eight GPUs, and only a few baselines achieve 99\% TPOT attainment when seven or eight GPUs are provisioned.

\begin{figure}[t]
    \centering
    \begin{subfigure}[b]{\linewidth}
        \centering
        \includegraphics[width=0.82\linewidth]{figs/microbench/microbench-rate-rep30-fixed.pdf}
    \end{subfigure}

    \begin{subfigure}[b]{\linewidth}
        \centering
        \includegraphics[width=0.82\linewidth]{figs/microbench/microbench-kvcache-rep30-fixed.pdf}
    \end{subfigure}

    \begin{subfigure}[b]{\linewidth}
        \centering
        \includegraphics[width=0.82\linewidth]{figs/microbench/microbench-tput-rep30-fixed.pdf}
    \end{subfigure}
   \vspace{-6mm}
    \caption{Benefits of cross-model memory coordination. The first figure shows the request rates. The last two figures shows the total KV memory size and the throughout of the two models.}
    \label{fig:elastic-mem-microbench}
\end{figure}

\mysubsection{Performance Analysis}
\label{sec:ablation}

Next, we provide a detailed performance breakdown to analyze the effectiveness of each design in \tool and how each contributes to its strong end-to-end performance.

\MyPara{Flexible cross-model memory sharing.}
We first evaluated the benefits of our flexible cross-model memory sharing by comparing with static partition using a simplified two-model trace extracted from Arena-Chat, shown in Figure~\ref{fig:elastic-mem-microbench} (first row).
We present the normalized total KV cache usage and aggregated throughput for both methods in the last two rows of Figure~\ref{fig:elastic-mem-microbench}.
As we can see, \tool's on-demand memory allocation allows it to use more memory for KV cache, particularly after the 20th second, when Model1 experiences low demand while Model2 faces a surge in request rates.
The larger KV cache memory enables \tool to achieve higher throughput, as shown in the last row. 
In contrast, under static partitioning, even when Model1 underutilizes its memory, Model2 cannot leverage the unused memory due to the static allocation boundary. 

\MyPara{Model placement.}
Next, we evaluated the benefits of our global scheduler that conducts model placement.
In this experiment, we used two GPUs to serve eight models from Arena-Chat. Figure~\ref{fig:global-scheduler-attainment} presents the TTFT and TPOT attainment with the global scheduler enabled or disabled.
The results show that enabling the global scheduler significantly improves both TTFT and TPOT attainment.
To provide further insights, we plot the average KV cache memory available per request as it arrives on each GPU. With the global scheduler enabled, the load is more evenly balanced across the two GPUs, allowing each request to access more KV cache memory on average.
In contrast, without the global placement scheduler, the load is imbalanced: GPU1 shows more available memory during the first 600 seconds, while the near-zero availability between the 800th and 1000th seconds indicates that GPU1 is idle while GPU0 is overloaded.

\begin{figure}[t]
  \centering
  \begin{subfigure}[t]{0.46\linewidth}
    \includegraphics[width=\linewidth]{figs/microbench/microbench-global-sch-attainment-fixed.pdf}
    \vspace{-6mm}
    \caption{Attainment with rate scales}
    \label{fig:global-scheduler-attainment}
  \end{subfigure}
  \begin{subfigure}[t]{0.49\linewidth}
    \includegraphics[width=\linewidth]{figs/microbench/microbench-global-sch-mem-per-req-fixed.pdf}
    \vspace{-6mm}
    \caption{GPU load status}
    \label{fig:global-scheduler-load}
  \end{subfigure}
  \caption{Effectiveness of global model placement scheduling.}
  \label{fig:global-scheduler}
\end{figure}

\begin{figure}[t]
  \centering
  \begin{subfigure}[t]{0.48\linewidth}
    \includegraphics[width=\linewidth]{figs/microbench/microbench-gpu-sch-slo-attainment-fixed.pdf}
    \caption{Attainment with SLO scales}
    \label{fig:gpu-scheduler-slo}
  \end{subfigure}
  \begin{subfigure}[t]{0.48\linewidth}
    \includegraphics[width=\linewidth]{figs/microbench/microbench-gpu-sch-queue_length-1-fixed.pdf}
    \caption{Queue length with time}
    \label{fig:gpu-scheduler-qlen}
  \end{subfigure}
  \caption{Effectiveness of GPU local request scheduling.}
  \label{fig:gpu-scheduler}
\end{figure}

\MyPara{Request arbitration.}
To evaluate the GPU-local scheduler that prioritizes requests, we use two models: we fix the SLO scale of Model1 to eight and vary the SLO scales of Model2 to evaluate the priority-based admission control in the GPU-local scheduling.
Figure~\ref{fig:gpu-scheduler-slo} shows the TTFT attainment as we vary the SLO scale of Model2. Model1 consistently maintains high attainment, and enabling our GPU-local scheduling improves the SLO attainment of Model2 by more than 40\%.
To dive deeper, we plot the queue length of each model in Figure~\ref{fig:gpu-scheduler-qlen} of one experiment run.
From the queue lengths, we clearly observe that when the local scheduler is enabled, the system prioritizes Model2's requests, which are shorter but have stricter SLO requirements. 

\mysubsection{Large-Scale Evaluation}
\label{sec:large-scale}

Finally, we evaluated how effectively \tool in reducing the cost of multi-LLM serving at scale.
We served all 58 models listed in Table~\ref{tab:models}, following common TP practices for large models: TP=4 for 32B models~\cite{qwen2.5, qwen2.5hf} and TP=4 or 8 for 70B models~\cite{llama3, llama3.3hf}, utilizing 32 GPUs in total.
We sampled 58 models from the Arena-Chat trace and the Hyperbolic trace.

\begin{figure}[t]
  \centering
  \begin{subfigure}[t]{0.98\linewidth}
    \includegraphics[width=\linewidth]{figs/slo_attainment_large_scale-2-fixed.pdf}
   \vspace{-6mm}
    \caption{SLO attainment with cluster size scaling}
    \label{fig:large-scale-attainment}
  \end{subfigure}
  \begin{subfigure}[t]{0.98\linewidth}
    \includegraphics[width=\linewidth]{figs/slo_attainment_gpu_need_99-fixed.pdf}
   \vspace{-4mm}
    \caption{Number of GPUs needed for 99\% SLO attainments}
    \label{fig:large-scale-gpus}
  \end{subfigure}
  \vspace{-2mm}
  \caption{SLO attainment and cost saving at large scales.}
  \label{fig:large-scale}
  \vspace{-2mm}
\end{figure}

\MyPara{SLO attainments vs. number of GPUs.} 
Figure~\ref{fig:large-scale-attainment} shows TTFT and TPOT attainment with an increased number of GPUs.
\tool outperforms all baselines, achieving nearly 99\% TTFT attainment with just 16 GPUs, while \muxp requires 32 GPUs to reach similar performance, and other methods require even more.
As the number of GPUs increases, all baselines improve in TTFT and TPOT attainment except QLM.
We find this to be related to its suboptimal scheduling algorithm.
QLM assigns requests to GPUs without considering which models are already on GPUs; when a queue is empty, it picks the first available GPU, often triggering unnecessary swaps. With more GPUs, this leads to more idle devices and frequent swapping. Its inefficient swapping, which requires engine restart~\cite{qlm2024github}, adds significant latency, causing queued requests to miss their SLOs.

\begin{figure}[t]
    \centering
    \includegraphics[width=0.90\linewidth]{figs/model_loading/activation_time_comparison_horizontal-fixed.pdf}
    \vspace{-3mm}
    \caption{The activation time for models with different sizes. Data is measured on H100 GPUs.}
    \vspace{-2mm}
    \label{fig:loading-time}
\end{figure}

\MyPara{Cost saving.} 
Figure~\ref{fig:large-scale-gpus} shows the number of GPUs required to achieve 99\% SLO attainments at different SLO scales.
If a system fails to achieve 99\% attainment with all 32 GPUs, we denote its GPU requirement as ``32+''. 
\tool achieves 99\% TTFT SLO attainments with only 16 GPUs when SLO scale is 5 and TPOT SLO scale is 2.0. 
\muxp needs 20 GPUs to get 99\% TTFT SLO attainments with SLO scale $\geq$ 30, while static partition needs even more GPUs or a higher SLO scale.
For TPOT, static partition is the best across all baselines, requiring 20 GPUs, while QLM and \muxp need at least 29 GPUs when TPOT SLO scale $\leq$ 3.0.

\begin{figure}[t]
  \centering
  \begin{subfigure}[t]{0.48\linewidth}
    \includegraphics[width=\linewidth]{figs/combined_gpu_utilization-fixed.pdf}
    \caption{Company A}
    \label{fig:production-throughput-a}
  \end{subfigure}
  \begin{subfigure}[t]{0.48\linewidth}
    \includegraphics[width=\linewidth]{figs/novita_revenue_comparison-fixed.pdf}
    \caption{Company B}
    \label{fig:production-throughput-b}
  \end{subfigure}
  \caption{Production results from two companies, showing (a) throughput and (b) revenue per GPU before and after using \tool.}
  \label{fig:production-throughput}
\end{figure}

\mysubsection{System Overhead}

\MyPara{Model activation latency.}
We also measured model activation latency from pageable CPU memory for models ranging from 1B to 70B parameters (see \S\ref{sec:model-loading} for more details).
Our optimizations significantly reduce activation latency. \tool loads small models (1B$\sim$8B) within 0.7s, a medium-size model (14B) in just 1.3s, and large models (>70B) in 1.5s (Figure~\ref{fig:loading-time}).
These results show that \tool can promptly activate evicted models upon receiving new requests.

\MyPara{Activation and migration frequency.}
Over a 10-minute window of the eight-models-on-two-GPUs run, \tool issued 2 idle-driven activations and 3 inter-GPU migrations. Migrations stay off the TTFT critical path: weights and in-flight KV state transfer over NVLink within tens of milliseconds (\eg, $\sim$20\,ms for an 8B model) while the source instance keeps serving. Activations complete within 2\,s, well inside the multi-second TTFT SLOs typical of online serving.

\MyPara{Elastic memory overhead.}
We evaluated \tool's elastic memory overhead in the worst case, where the request rates are constant, leaving no opportunities for dynamic sharing. With two colocated Llama-3.2-3B models on an A100-40G GPU, \tool incurs
 only 3 ms (4 \%) TTFT and 4 ms (13 \%) TPOT overhead at a high load of 32 req/s and
only 2 ms (3.5 \%) TTFT and 3 ms (7 \%) TPOT overhead at a high load of 28 req/s (see~\S\ref{sec:elastic-overhead}), compared to static partitioning.

\mysubsection{Production Workloads}
As of December 2025, \tool has been deployed in multiple organizations serving LLM workloads. Figure~\ref{fig:production-throughput} summarizes deployments at a large tech company (Company A) using \tool on an internal GPU cluster, and a commercial LLM inference provider (Company B) serving external customers.
To isolate \tool's contribution from workload variation, both deployments use shadow workload replay: the same online request stream is mirrored to two identical clusters, one running \tool and the other not, so that the model mix and traffic patterns are held constant across the two arms.

Figure~\ref{fig:production-throughput-a} shows Company A’s per-GPU token throughput before and after adopting \tool. Company A serves diverse applications using fine-tuned and off-the-shelf models (3B–70B) on a shared GPU pool. With \tool dynamically adjusting model placement and sharing\textemdash without changes to existing inference engines\textemdash Company A achieves consistently higher throughput (3.89× on average) over several weeks, with no SLO violations and unchanged tail latency. Figure~\ref{fig:production-throughput-b} presents results from Company B, which optimizes revenue per GPU under highly variable traffic, computed as the prefill and decode tokens generated in the window priced at the provider's published per-token rates, normalized by GPU count. Prior to \tool, fragmentation and overprovisioning constrained revenue; after deployment, \tool improves utilization while meeting all SLOs, raising revenue per GPU by 2.86$\times$ by converting idle capacity into billable tokens.

